% ============================================================
\chapter{Introduction to NLP}
\label{ch:introduction}
% ============================================================

In 1950, Alan Turing posed the founding question of artificial intelligence:
can a machine ``think''? And he proposed an operational criterion: if a
human, conversing with a machine via text, cannot tell it apart from another
human, then the machine ``thinks'' in a functional sense. This ``Turing
test'' is fundamentally a \emph{natural language processing} challenge: to
pass it, a machine must understand syntax, semantics, context, irony, and
subtext. Seventy-five years later, large language models (GPT, Claude, LLaMA)
seem to be approaching this threshold --- yet the fundamental challenges
remain intact.

\begin{intuition}[title=\textbf{Why study NLP?}]
Language is the primary vehicle for human communication. Teaching machines to
understand, produce, and reason about natural language is one of the fundamental
challenges of artificial intelligence. Applications range from machine
translation to information extraction, conversational agents, and text
generation.
\end{intuition}

% ------------------------------------------------------------
\section{What is NLP?}
\label{sec:intro:what}
% ------------------------------------------------------------

\begin{definition}[Natural Language Processing]
\emph{Natural Language Processing} (NLP) is the subfield of artificial
intelligence that studies interactions between computers and human language. It
encompasses the analysis, understanding, and generation of text and speech in
natural languages.
\end{definition}

NLP is distinguished from \emph{computational linguistics} by its
application-oriented focus, although both fields share a common theoretical
foundation. Formally, we can view language as a generative system defined over
a vocabulary $\mathcal{V}$:
\[
\mathcal{L} \subseteq \mathcal{V}^* = \bigcup_{n=0}^{\infty} \mathcal{V}^n
\]
where $\mathcal{V}^*$ denotes the set of all finite sequences over $\mathcal{V}$.

% ------------------------------------------------------------
\section{Levels of Linguistic Analysis}
\label{sec:intro:levels}
% ------------------------------------------------------------

Natural language processing operates at several levels of analysis:

\begin{enumerate}
    \item \textbf{Phonetics and phonology}: study of speech sounds (primarily
          for speech processing).
    \item \textbf{Morphology}: internal structure of words (prefixes, suffixes,
          inflections). For example, ``un-believe-able'' decomposes into
          \texttt{un-} + \texttt{believe} + \texttt{-able}.
    \item \textbf{Syntax}: sentence structure, grammatical relations.
    \item \textbf{Semantics}: meaning of words and sentences.
    \item \textbf{Pragmatics}: meaning in context, speaker intentions.
    \item \textbf{Discourse}: coherence and structure beyond the sentence.
\end{enumerate}

\begin{remark}
In modern NLP, deep models implicitly learn these levels of analysis from data,
without explicit segmentation. However, understanding these levels remains
essential for diagnosing errors and designing relevant evaluations.
\end{remark}

% ------------------------------------------------------------
\section{Fundamental NLP Tasks}
\label{sec:intro:tasks}
% ------------------------------------------------------------

\subsection{Text Classification}

Text classification assigns a label $y \in \mathcal{Y}$ to a document
$\mathbf{x} = (x_1, x_2, \ldots, x_T)$. The model learns a function
$f : \mathcal{V}^* \to \mathcal{Y}$ from a training set
$\{(\mathbf{x}_i, y_i)\}_{i=1}^{N}$.

Applications: sentiment analysis, spam detection, topic classification.

\subsection{Sequence Labeling}

Sequence labeling assigns a label $y_t \in \mathcal{Y}$ to each token $x_t$
in a sequence:
\[
f : \mathcal{V}^T \to \mathcal{Y}^T
\]

Applications: named entity recognition (NER), part-of-speech (POS) tagging.

\begin{example}
Consider the sentence ``Mary lives in Paris'' with the following NER labeling:
\begin{center}
\begin{tabular}{cccc}
Mary & lives & in & Paris \\
\texttt{B-PER} & \texttt{O} & \texttt{O} & \texttt{B-LOC}
\end{tabular}
\end{center}
\end{example}

\subsection{Machine Translation}

Machine translation is a sequence-to-sequence transduction task:
\[
f : \mathcal{V}_{\text{source}}^* \to \mathcal{V}_{\text{target}}^*
\]
It is historically one of the driving tasks of NLP.

\subsection{Question Answering}

Given a context $C$ and a question $Q$, the system must extract or generate an
answer $A$:
\[
f : (C, Q) \to A
\]

\subsection{Automatic Summarization}

Automatic summarization produces a condensed version $S$ of a document $D$,
preserving essential information: $f : D \to S$ with $\abs{S} \ll \abs{D}$.

% ------------------------------------------------------------
\section{Probabilistic Foundations}
\label{sec:intro:proba}
% ------------------------------------------------------------

Probabilistic language modeling is at the heart of NLP. A \emph{language model}
defines a probability distribution over token sequences.

\begin{definition}[Language Model]
A \emph{language model} is a probability distribution $P$ over $\mathcal{V}^*$
such that:
\[
P(\mathbf{x}) = P(x_1, x_2, \ldots, x_T) = \prod_{t=1}^{T} P(x_t \mid x_1, \ldots, x_{t-1})
\]
where the factorization follows from the chain rule of probability.
\end{definition}

\begin{theorem}[Chain Rule]
For any joint distribution $P(x_1, \ldots, x_T)$:
\[
P(x_1, \ldots, x_T) = \prod_{t=1}^{T} P(x_t \mid x_{<t})
\]
where $x_{<t} = (x_1, \ldots, x_{t-1})$ denotes the \emph{context} at step $t$.
\end{theorem}

\begin{proof}
By successive applications of the definition of conditional probability
$P(A \mid B) = P(A, B) / P(B)$:
\begin{align*}
P(x_1, x_2, \ldots, x_T) &= P(x_1) \cdot P(x_2 \mid x_1) \cdot P(x_3 \mid x_1, x_2) \cdots \\
                           &= \prod_{t=1}^{T} P(x_t \mid x_{<t}). \qedhere
\end{align*}
\end{proof}

% ------------------------------------------------------------
\section{Information Theory and Language}
\label{sec:intro:info}
% ------------------------------------------------------------

\begin{definition}[Entropy]
The \emph{entropy} of a discrete random variable $X$ taking values in
$\mathcal{X}$ is:
\[
H(X) = -\sum_{x \in \mathcal{X}} P(x) \log_2 P(x)
\]
It measures the average uncertainty associated with $X$.
\end{definition}

\begin{definition}[Cross-Entropy]
The \emph{cross-entropy} between the true distribution $p$ and a model $q$ is:
\[
H(p, q) = -\sum_{x} p(x) \log q(x) = H(p) + \KL(p \| q)
\]
\end{definition}

\begin{proposition}[Lower Bound]
For any model $q$, we have $H(p, q) \geq H(p)$, with equality if and only if
$q = p$.
\end{proposition}

The \emph{perplexity} of a language model $q$ is defined as:
\[
\text{PPL}(q) = 2^{H(p, q)}
\]

\begin{keyformulas}[title=\textbf{Key Information-Theoretic Metrics}]
\begin{align}
H(X) &= -\sum_{x} P(x) \log P(x) & \text{(entropy)} \\
\KL(p \| q) &= \sum_{x} p(x) \log \frac{p(x)}{q(x)} & \text{(KL divergence)} \\
\text{PPL} &= \exp\!\left(-\frac{1}{T}\sum_{t=1}^{T} \log q(x_t \mid x_{<t})\right) & \text{(perplexity)}
\end{align}
\end{keyformulas}

% ------------------------------------------------------------
\section{Classical NLP Pipeline}
\label{sec:intro:pipeline}
% ------------------------------------------------------------

A classical NLP pipeline comprises the following stages:

\begin{enumerate}
    \item \textbf{Data collection}: text corpora, web scraping, APIs.
    \item \textbf{Preprocessing}:
    \begin{itemize}
        \item Tokenization (splitting into lexical units)
        \item Normalization (lowercasing, accent removal)
        \item Stopword removal
        \item Stemming or lemmatization
    \end{itemize}
    \item \textbf{Representation}: converting text to numerical vectors.
    \item \textbf{Modeling}: training a model.
    \item \textbf{Evaluation}: measuring performance on a test set.
    \item \textbf{Deployment}: putting the model into production.
\end{enumerate}

\begin{warning}[title=\textbf{Preprocessing and modern models}]
With pre-trained models (BERT, GPT), preprocessing is generally minimal: the
model's tokenizer (BPE, WordPiece) replaces classical stemming and stopword
removal steps. Applying aggressive preprocessing can actually degrade
performance.
\end{warning}

% ------------------------------------------------------------
\section{Tokenization}
\label{sec:intro:tokenization}
% ------------------------------------------------------------

\begin{definition}[Tokenization]
\emph{Tokenization} is the process of splitting text into elementary units
called \emph{tokens}. A token can be a word, a subword, a character, or a byte.
\end{definition}

\subsection{Subword Tokenization}

Modern subword tokenization algorithms (BPE, WordPiece, Unigram) offer a
trade-off between word tokenization (very large vocabulary) and character
tokenization (very long sequences).

\begin{algorithme}[title=\textbf{Byte Pair Encoding (BPE)}]
\begin{enumerate}
    \item Initialize the vocabulary with all individual characters.
    \item Repeat $k$ times:
    \begin{enumerate}
        \item Count all adjacent symbol pairs in the corpus.
        \item Merge the most frequent pair into a new symbol.
        \item Add this symbol to the vocabulary.
    \end{enumerate}
    \item Return the final vocabulary of size $\abs{\mathcal{V}_0} + k$.
\end{enumerate}
\end{algorithme}

\begin{codeblock}[title=\textbf{Tokenization with Hugging Face}]
\begin{minted}{python}
from transformers import AutoTokenizer

tokenizer = AutoTokenizer.from_pretrained("bert-base-uncased")
text = "Natural language processing is fascinating."
tokens = tokenizer.tokenize(text)
print(tokens)
# ['natural', 'language', 'processing', 'is', 'fascinating', '.']

ids = tokenizer.encode(text, return_tensors="pt")
print(ids)
# tensor([[ 101, 3019, 2653, 6364, 2003, 15281, 1012,  102]])
\end{minted}
\end{codeblock}

% ------------------------------------------------------------
\section{Brief History of NLP}
\label{sec:intro:history}
% ------------------------------------------------------------

\begin{center}
\renewcommand{\arraystretch}{1.2}
\begin{tabular}{lll}
\toprule
\textbf{Period} & \textbf{Paradigm} & \textbf{Examples} \\
\midrule
1950--1970 & Symbolic rules & ELIZA, SHRDLU \\
1970--1990 & Formal grammars & ATN, LFG, HPSG \\
1990--2010 & Statistical methods & HMM, CRF, $n$-grams \\
2010--2017 & Deep learning & Word2Vec, LSTM, Seq2Seq \\
2017--      & Transformers \& LLMs & BERT, GPT, T5, LLaMA \\
\bottomrule
\end{tabular}
\end{center}

\begin{remark}
The transition from the statistical paradigm to the neural paradigm does not
represent a complete break: the concepts of cross-entropy, language modeling,
and regularization remain at the heart of modern approaches.
\end{remark}

% ------------------------------------------------------------
\section{Open Challenges}
\label{sec:intro:challenges}
% ------------------------------------------------------------

Despite spectacular progress, many challenges remain open:

\begin{itemize}
    \item \textbf{Real understanding vs.\ statistical correlations}: do LLMs
          capture \emph{meaning} or merely surface regularities?
    \item \textbf{Low-resource languages}: the majority of resources are
          available in English; approximately 7,000 languages are spoken
          worldwide.
    \item \textbf{Robustness}: sensitivity to adversarial perturbations,
          neologisms, out-of-distribution domains.
    \item \textbf{Efficiency}: LLMs require considerable computational
          resources.
    \item \textbf{Ethics}: bias, toxicity, privacy, misinformation.
\end{itemize}

% ------------------------------------------------------------
\section{Exercises}
\label{sec:intro:exercises}
% ------------------------------------------------------------

\begin{exercise}[$\star$]
Compute the entropy of a uniform language model over a vocabulary of size
$V = 50{,}000$. What is the corresponding perplexity?
\end{exercise}

\begin{exercise}[$\star$]
Consider the sentence ``The cat eats the mouse.'' Apply the chain rule to
express $P(\text{The, cat, eats, the, mouse})$ as a product of conditional
probabilities.
\end{exercise}

\begin{exercise}[$\star\star$]
Implement a simple BPE tokenizer in Python. Start with a small corpus of 10
sentences and show the vocabulary evolution after 20 merges.
\end{exercise}

\begin{exercise}[$\star\star$]
Show that the Kullback-Leibler divergence $\KL(p \| q)$ is always non-negative
(Gibbs' inequality). \emph{Hint:} use the inequality $\log x \leq x - 1$.
\end{exercise}

\begin{exercise}[$\star\star\star$]
Compare the perplexity of an $n$-gram model (for $n = 1, 2, 3$) and a neural
model (GPT-2) on a corpus of your choice. Analyze the results as a function of
training corpus size.
\end{exercise}
